\section{Discussion}
\noindent\textbf{Multimodal Understanding.}\quad
While text is capable of conveying an enormous number of concepts about the world, many important concepts are conveyed mainly through other modalities, such as images, audio, and physical interaction \citep{bisk2020experiencegroundslang}. Existing large-scale NLP models, such as GPT-3, do not incorporate multimodal information, so we design our benchmark to capture a diverse array of tasks in a text-only format. However, as models gain the ability to process multimodal inputs, benchmarks should be designed to reflect this change. One such benchmark could be a ``Turk Test,'' consisting of Amazon Mechanical Turk Human Intelligence Tasks. These are well-defined tasks that require models to interact with flexible formats and demonstrate multimodal understanding.

\noindent\textbf{The Internet as a Training Set.}\quad
A major distinction between our benchmark and previous multitask NLP benchmarks is that we do not require large training sets. Instead, we assume that models have acquired the requisite knowledge from reading vast quantities of diverse text from the Internet. This process is typically called pretraining, but it can be thought of as training in its own right, where the downstream evaluation is demonstrating whatever knowledge we would expect a human to pick up from reading the same text. 
% This methodological change enables collecting a much more extensive and diverse set of tasks for evaluation, because we do not need large training sets that exactly match the evaluation format for each task.

%These factors motivate the design of our benchmark, for which we use a new methodology without large IID training sets. Instead, we have a dev, val, and test set for each task. The dev set is used for few-shot prompts, the val set could be used for hyperparameter tuning, and the test set is used to compute the final accuracy. Importantly, the format of our evaluation is not identical to the format in which information is acquired at training time. This has the benefit of drastically reducing spurious cues [cite adversarial filtering] and is in stark contrast to the previous paradigm of identically formatted training and test sets. We anticipate our methodology becoming more widespread as models improve at extracting information from diverse online sources.

This motivates us to propose a methodological change so that models are trained more like how humans learn.
While most previous machine learning benchmarks have models learn from a large question bank, humans primarily learn new subjects by reading books and listening to others talk about the topic. For specialized subjects such as Professional Law, massive legal corpora are available, such as the 164-volume legal encyclopedia \emph{Corpus Juris Secundum}, but there are fewer than 5,000 multistate bar exam questions available. Learning the entire law exclusively through a small number of practice tests is implausible, so future models must learn more during pretraining. 

For this reason we assess pretrained models in a zero-shot, few-shot, or transfer setting and we provide a dev, val, and test set for each task. The dev set is used for few-shot prompts, the val set could be used for hyperparameter tuning, and the test set is used to compute the final accuracy. Importantly, the format of our evaluation is not identical to the format in which information is acquired during pretraining. This has the benefit of obviating concerns about spurious training set annotation artifacts \citep{geirhos2020shortcut,Hendrycks2019NaturalAE} and is in stark contrast to the previous paradigm of identically distributed training and test sets. 
This change also enables collecting a much more extensive and diverse set of tasks for evaluation.
We anticipate our methodology becoming more widespread as models improve at extracting information from diverse online sources.

%Without the need for immense training data identically distributed to the test data, we can assess model performance across a diverse set of subjects and more easily find blindspots in the pretraining process. At the same time, this design choice also obviates concerns of spurious training set annotation artifacts [], so our test provides a more reliable estimate of accuracy.

%While we find that large-scale Transformers have surprisingly high performance on many subjects, none of the models we test excel in any individual field, including at the introductory level. This greatly differs from human learning, where many subjects can be mastered at the introductory level early on, and a few can be mastered at the professional level throughout the course of a lifetime. Moreover, the models we test perform poorly on subject areas that require numerical calculations, such as elementary mathematics.

\noindent\textbf{Model Limitations.}\quad
We find that current large-scale Transformers have wide room for improvement. They are notably poor at modeling human (dis)approval, as evident by the low performance on the Professional Law and Moral Scenarios tasks. For future systems to be aligned with human values, high performance on these tasks is crucial \citep{hendrycks2020ethicsdataset}, so future research should especially aim to increase accuracy on these tasks. Models also have difficulty performing calculations, so much so that they exhibit poor performance on Elementary Mathematics and many other STEM subjects with ``plug and chug'' problems. Additionally, they do not match expert-level performance (90\%) on any subject, so for all subjects it is subhuman. On average, models are only now starting to move beyond random-chance accuracy levels.

Addressing these shortcomings may be challenging. To illustrate this, we attempted to create a better Professional Law model by pretraining on specialized data but achieved only limited success. We collected approximately 2,000 additional Professional Law training examples. After fine-tuning a RoBERTa-base model \citep{RobertaLiu2019AR} using this custom training set, our model attained $32.8\%$ test accuracy. To test the impact of additional specialized training data, we also had RoBERTa continue pretraining on approximately 1.6 million legal case summaries using Harvard’s Law Library case law corpus \texttt{case.law}, but after fine-tuning it only attained $36.1\%$ accuracy. This suggests that while additional pretraining on relevant high quality text can help, it may not be enough to substantially increase the performance of current models. 

It is unclear whether simply scaling up existing language models will solve the test. Current understanding indicates that a $10\times$ increase in model size must be accompanied by an approximate $5\times$ increase in data \citep{kaplan2020scalinglaws}. Aside from the tremendous expense in creating multi-trillion parameter language models, data may also become a bottleneck, as there is far less written about esoteric branches of knowledge than about everyday situations.

% we attempt to create a better model for Professional Law. 
% We speculate that data efficiency during pretraining may pose a challenge toward better.
% We speculate that models must be more efficient during the language modeling pretraining phase.
% assess
% We tested whether  with 

% Inefficient:
% In this section we compare fine-tuned performance with and without additional pretraining. We focus on Professional Law, for which we collected an additional $1$k training examples for fine-tuning and $10$GB of extra law-related pretraining documents. We compare RoBERTa-base [] after fine-tuning and after pretraining on the additional $10$GB of law documents before fine-tuning. We found that few-shot GPT-3 gets $33.9\%$ accuracy, while fine-tuned RoBERTa does slightly worse at $32.8\%$. RoBERTa with extra pretraining on law documents did only slightly better, at $36.1\%$ accuracy. This suggests that while additional pretraining on relevant high quality text can help, it may not be enough to substantially increase the performance of current models.
% More efficient.
% It is unclear whether simply scaling up existing language models will solve these tasks [cite GPT-3]. Current understanding indicates that a $10\times$ increase in model size must be accompanied by an approximate $5\times$ increase in data [cite scaling laws]. This may become a bottleneck for existing models, because there is far less data available for specialized branches of knowledge than for everyday topics.

% maybe let's add some pro-safety stuff in the discussion